
\documentclass[11pt,a4paper]{article}

\usepackage[utf8]{inputenc}
\usepackage[T1]{fontenc}
\usepackage[ngerman]{babel}
\usepackage{amsmath,amssymb}
\usepackage{geometry}
\geometry{margin=2.5cm}

\title{EABC-Holonomien, Primvierlinge und die Riemannsche Zetafunktion}
\author{Thomas Hoffbauer}
\date{\today}

\begin{document}

\maketitle

\begin{abstract}
Wir untersuchen ein experimentelles Modell zur Beschreibung von Primzahlvierlingen
mittels eines zyklischen EABC-Phasenfeldes. Numerische Tests zeigen,
dass die Bandstruktur des Feldes signifikante Korrelationen mit
kosinusartigen Komponenten eines aus den ersten Nullstellen der
Riemannschen Zetafunktion konstruierten Spektrums besitzt.
\end{abstract}

\section{Historischer Hintergrund}

Euler erkannte die Verbindung zwischen Primzahlen und der Zetafunktion.
Riemann zeigte später, dass die Nullstellen der Zetafunktion
mit den Fluktuationen der Primzahldichte zusammenhängen.

\section{Das EABC-Modell}

Primzahlen modulo 12 werden in vier Kanäle zerlegt:

\[
E \equiv 1,\;
A \equiv 5,\;
B \equiv 7,\;
C \equiv 11 \pmod{12}.
\]

Primvierlinge erzeugen zyklische Holonomiestrukturen.

\section{Nullstellentest}

Die stärkste Korrelation entsteht zwischen der Bandenergie
des Feldes und kosinusartigen Komponenten der Zetafunktion:

\[
\mathrm{corr}(|\Psi|,Z_{\cos}) \approx -0.24
\]

für die ersten 2000 Nullstellen.

\section{Interpretation}

Die Ergebnisse sprechen dafür, dass die Zetafunktion vor allem
globale Dichteinformationen der Primzahlen erfasst.

\end{document}
